\documentclass[12pt,a4paper]{article}

\usepackage[utf8]{inputenc}
\usepackage[T5]{fontenc}
\usepackage[vietnamese]{babel}
\usepackage{amsmath}
\usepackage{geometry}
\usepackage{listings}
\usepackage{xcolor}

\geometry{
    a4paper,
    left=2.5cm,
    right=2.5cm,
    top=2.5cm,
    bottom=2.5cm
}

\lstset{
    basicstyle=\ttfamily\small,
    breaklines=true,
    breakatwhitespace=false,
    columns=flexible,
    keepspaces=true,
    frame=single
}

\begin{document}

\section{Practice 1: GPIO Control of LED, Button, Relay and Buzzer}

\subsection{Objective}

The objective of this practice is to use the GPIO pins of the ATmega328P
to read a push button and control an LED, a relay, and a buzzer.

The button is configured as a digital input, while the LED, relay, and
buzzer are configured as digital outputs.

\subsection{Pin Configuration}

The pins used in this practice are defined as follows:

\begin{lstlisting}[language=C]
const int BUTTON_PIN = 2;
const int LED_PIN = 13;
const int RELAY_PIN = 8;
const int BUZZER_PIN = 9;
\end{lstlisting}

The corresponding GPIO functions are:

\begin{center}
\begin{tabular}{|c|c|c|}
\hline
\textbf{Pin} & \textbf{Device} & \textbf{Function} \\
\hline
D2 & Button & Digital input \\
\hline
D13 & LED & Digital output \\
\hline
D8 & Relay & Digital output \\
\hline
D9 & Buzzer & Digital output \\
\hline
\end{tabular}
\end{center}

\subsection{GPIO Initialization}

The button is configured as an input, while the LED, relay, and buzzer
are configured as outputs.

\begin{lstlisting}[language=C]
pinMode(BUTTON_PIN, INPUT);
pinMode(LED_PIN, OUTPUT);
pinMode(RELAY_PIN, OUTPUT);
pinMode(BUZZER_PIN, OUTPUT);
\end{lstlisting}

The initial output states are set so that all devices are turned off:

\begin{lstlisting}[language=C]
digitalWrite(LED_PIN, LOW);
digitalWrite(RELAY_PIN, HIGH);
digitalWrite(BUZZER_PIN, LOW);
\end{lstlisting}

The LED and buzzer are active-high, while the relay is active-low.

Therefore:

\[
\begin{array}{c|c|c}
\textbf{Device} & \textbf{ON} & \textbf{OFF} \\
\hline
LED & HIGH & LOW \\
Relay & LOW & HIGH \\
Buzzer & HIGH & LOW
\end{array}
\]

\subsection{Reading the Button}

The state of the button is read using the digital input function:

\begin{lstlisting}[language=C]
int buttonState = digitalRead(BUTTON_PIN);
\end{lstlisting}

The button uses active-low logic. Therefore, a LOW signal indicates
that the button is pressed:

\[
\begin{array}{c|c}
\textbf{Button State} & \textbf{Meaning} \\
\hline
LOW & Pressed \\
HIGH & Released
\end{array}
\]

The program checks whether the button state is LOW:

\begin{lstlisting}[language=C]
if (buttonState == LOW)
{
    ...
}
else
{
    ...
}
\end{lstlisting}

\subsection{Controlling the LED, Relay and Buzzer}

When the button is pressed, the LED and buzzer are turned on by
outputting HIGH, while the relay is turned on by outputting LOW.

\begin{lstlisting}[language=C]
digitalWrite(LED_PIN, HIGH);
digitalWrite(RELAY_PIN, LOW);
digitalWrite(BUZZER_PIN, HIGH);
\end{lstlisting}

When the button is released, all devices are turned off:

\begin{lstlisting}[language=C]
digitalWrite(LED_PIN, LOW);
digitalWrite(RELAY_PIN, HIGH);
digitalWrite(BUZZER_PIN, LOW);
\end{lstlisting}

The complete control logic is:

\[
\text{Button LOW (Pressed)}
\rightarrow
\begin{cases}
\text{LED HIGH} & \rightarrow \text{LED ON}\\
\text{Relay LOW} & \rightarrow \text{Relay ON}\\
\text{Buzzer HIGH} & \rightarrow \text{Buzzer ON}
\end{cases}
\]

When the button is released:

\[
\text{Button HIGH (Released)}
\rightarrow
\begin{cases}
\text{LED LOW} & \rightarrow \text{LED OFF}\\
\text{Relay HIGH} & \rightarrow \text{Relay OFF}\\
\text{Buzzer LOW} & \rightarrow \text{Buzzer OFF}
\end{cases}
\]

\subsection{Main Program}

The program continuously reads the button state and updates the three
output devices according to the button state.

\begin{lstlisting}[language=C]
void loop()
{
    int buttonState = digitalRead(BUTTON_PIN);

    if (buttonState == LOW)
    {
        digitalWrite(LED_PIN, HIGH);
        digitalWrite(RELAY_PIN, LOW);
        digitalWrite(BUZZER_PIN, HIGH);
    }
    else
    {
        digitalWrite(LED_PIN, LOW);
        digitalWrite(RELAY_PIN, HIGH);
        digitalWrite(BUZZER_PIN, LOW);
    }
}
\end{lstlisting}

\subsection{Result}

The ATmega328P successfully reads the digital state of the push button
through a GPIO input and controls the LED, relay, and buzzer.

When the button is pressed, the LED is turned on, the relay is
activated, and the buzzer is turned on.

When the button is released, the LED, relay, and buzzer are turned off.

This practice demonstrates the basic operation of GPIO input and
output on the ATmega328P, including the use of active-low input and
active-low output logic.


\section{Practice 2: ADC and Sensor}

\subsection{Objective}

The objective of this practice is to use the ATmega328 ADC to read an analog signal, collect 16 samples, calculate the average ADC value, convert the result into voltage, and send the result through serial communication.

\subsection{ADC Configuration}

The ADC is configured with the following parameters:

\begin{itemize}
    \item ADC resolution: 10-bit
    \item ADC input channel: ADC2 (PC2)
    \item ADC reference voltage: AVCC = 5 V
    \item ADC prescaler: 128
    \item ADC conversion method: ADC interrupt
    \item Number of samples: 16
\end{itemize}

The ADC reference voltage is selected using the \texttt{REFS0} bit in the \texttt{ADMUX} register.

\[
V_{\mathrm{REF}} = V_{\mathrm{CC}} = 5\,\mathrm{V}
\]

ADC2 is selected by setting the \texttt{MUX1} bit.

The ADC prescaler is set to 128 using the \texttt{ADPS2}, \texttt{ADPS1}, and \texttt{ADPS0} bits.

For a CPU clock of 16 MHz, the ADC clock is:

\[
f_{\mathrm{ADC}}
=
\frac{f_{\mathrm{CPU}}}{128}
=
\frac{16\,\mathrm{MHz}}{128}
=
125\,\mathrm{kHz}
\]

\subsection{ADC Initialization}

The ADC is initialized by configuring the \texttt{ADMUX} and \texttt{ADCSRA} registers.

\begin{lstlisting}[language=C]
void ADC_Init() {
    ADMUX |= (1 << REFS0) | (1 << MUX1);

    ADCSRA = (1 << ADEN) | (1 << ADIE)
           | (1 << ADPS2) | (1 << ADPS1)
           | (1 << ADPS0);
}
\end{lstlisting}

The \texttt{ADEN} bit enables the ADC, while the \texttt{ADIE} bit enables the ADC conversion-complete interrupt.

\subsection{ADC Interrupt}

When an ADC conversion is completed, the ADC interrupt service routine is executed automatically.

\begin{lstlisting}[language=C]
ISR(ADC_vect) {
    adc_sum += ADC;
    sample_count++;

    if (sample_count >= 16) {
        adc_ready = true;
    } else {
        ADCSRA |= (1 << ADSC);
    }
}
\end{lstlisting}

The 10-bit ADC result is added to \texttt{adc\_sum}. After 16 conversions, the program sets the \texttt{adc\_ready} flag.

If fewer than 16 samples have been collected, the next ADC conversion is started by setting the \texttt{ADSC} bit.

\subsection{Sampling and Averaging}

The program collects 16 ADC samples before processing the data.

Let the ADC results be:

\[
ADC_1, ADC_2, \ldots, ADC_{16}
\]

The average ADC value is calculated as:

\[
ADC_{\mathrm{avg}}
=
\frac{1}{16}
\sum_{i=1}^{16} ADC_i
\]

The program calculates the average using:

\begin{lstlisting}[language=C]
uint16_t avg_adc = adc_sum / 16;
\end{lstlisting}

Averaging multiple samples provides a more stable ADC value by reducing small variations between individual readings.

\subsection{Conversion from ADC Value to Voltage}

The ATmega328 has a 10-bit ADC, so the digital output ranges from 0 to 1023.

The input voltage is calculated using:

\[
V_{\mathrm{IN}}
=
\frac{ADC_{\mathrm{avg}} \times V_{\mathrm{REF}}}{1024}
\]

With:

\[
V_{\mathrm{REF}} = 5\,\mathrm{V}
\]

the equation becomes:

\[
V_{\mathrm{IN}}
=
\frac{ADC_{\mathrm{avg}} \times 5}{1024}
\]

The program implements this conversion as:

\begin{lstlisting}[language=C]
float voltage = (avg_adc * 5.0) / 1024.0;
\end{lstlisting}

\subsection{Serial Output}

The averaged ADC value and the corresponding voltage are sent to the computer through UART at 9600 baud.

\begin{lstlisting}[language=C]
Serial.begin(9600);

Serial.print("Raw ADC: ");
Serial.print(avg_adc);
Serial.print(" | Voltage: ");
Serial.print(voltage);
Serial.println(" V");
\end{lstlisting}

The output has the following format:

\begin{lstlisting}
Raw ADC: <ADC value> | Voltage: <voltage> V
\end{lstlisting}

\subsection{Complete Program}

\begin{lstlisting}[language=C]
#include <Arduino.h>
#include <avr/interrupt.h>

volatile uint16_t adc_sum = 0;
volatile uint8_t sample_count = 0;
volatile bool adc_ready = false;

void ADC_Init() {
    ADMUX |= (1 << REFS0) | (1 << MUX1);

    ADCSRA = (1 << ADEN) | (1 << ADIE)
           | (1 << ADPS2) | (1 << ADPS1)
           | (1 << ADPS0);
}

ISR(ADC_vect) {
    adc_sum += ADC;
    sample_count++;

    if (sample_count >= 16) {
        adc_ready = true;
    } else {
        ADCSRA |= (1 << ADSC);
    }
}

void setup() {
    Serial.begin(9600);
    ADC_Init();
}

void loop() {
    if (!adc_ready && sample_count == 0) {
        ADCSRA |= (1 << ADSC);
    }

    if (adc_ready) {
        uint16_t avg_adc = adc_sum / 16;

        float voltage = (avg_adc * 5.0) / 1024.0;

        Serial.print("Raw ADC: ");
        Serial.print(avg_adc);
        Serial.print(" | Voltage: ");
        Serial.print(voltage);
        Serial.println(" V");

        adc_sum = 0;
        sample_count = 0;
        adc_ready = false;

        delay(500);
    }
}

\end{lstlisting}

\section{Practice 3: PWM and Control Servo Motor}

\subsection{Objective}

The objective of this practice is to generate PWM signals to control a
servo motor and an LED using the ATmega328. The PWM signal of the servo
motor is also observed using an oscilloscope.

\subsection{Servo Motor PWM}

The servo motor is connected to pin PB1, which corresponds to the OC1A
output of Timer1.

Timer1 is configured in Fast PWM Mode 14. The prescaler is set to 8 and
the value of \texttt{ICR1} is used to determine the PWM period.

\begin{lstlisting}[language=C]
void Servo_PWM_Init(void) {
    DDRB |= (1 << SERVO_PIN);

    TCCR1A = (1 << COM1A1) | (1 << WGM11);

    TCCR1B = (1 << WGM13) |
             (1 << WGM12) |
             (1 << CS11);

    ICR1 = 39999;
}
\end{lstlisting}

With a CPU clock of 16 MHz and a prescaler of 8, the Timer1 clock is:

\[
f_{\mathrm{Timer1}}
=
\frac{16\,\mathrm{MHz}}{8}
=
2\,\mathrm{MHz}
\]

Therefore, one timer count corresponds to:

\[
T_{\mathrm{count}}
=
\frac{1}{2\,\mathrm{MHz}}
=
0.5\,\mu\mathrm{s}
\]

The PWM period is:

\[
T_{\mathrm{PWM}}
=
(ICR1+1)T_{\mathrm{count}}
\]

With \texttt{ICR1 = 39999}:

\[
T_{\mathrm{PWM}}
=
40000 \times 0.5\,\mu\mathrm{s}
=
20\,\mathrm{ms}
\]

Therefore, the PWM frequency is:

\[
f_{\mathrm{PWM}}
=
\frac{1}{20\,\mathrm{ms}}
=
50\,\mathrm{Hz}
\]

Thus, Timer1 generates a PWM signal with a frequency of approximately
50 Hz for the servo motor.

\subsection{Servo Angle Control}

The servo angle is controlled by changing the pulse width of the PWM
signal through the \texttt{OCR1A} register.

The program limits the angle to the range from 0 to 180 degrees.

\begin{lstlisting}[language=C]
void Servo_SetAngle(uint8_t angle) {
    if (angle > 180)
        angle = 180;

    uint16_t min_ocr = 999;
    uint16_t max_ocr = 4999;

    OCR1A = min_ocr +
            ((uint32_t)angle *
            (max_ocr - min_ocr)) / 180;
}
\end{lstlisting}

The pulse width is approximately mapped from 0.5 ms to 2.5 ms:

\[
PW(\theta)
=
0.5\,\mathrm{ms}
+
\frac{\theta}{180}
(2.5-0.5)\,\mathrm{ms}
\]

where:

\[
0^\circ \leq \theta \leq 180^\circ
\]

Therefore:

\[
\theta = 0^\circ
\quad\Rightarrow\quad
PW \approx 0.5\,\mathrm{ms}
\]

and:

\[
\theta = 180^\circ
\quad\Rightarrow\quad
PW \approx 2.5\,\mathrm{ms}
\]

Changing the pulse width changes the position of the servo motor.

\subsection{LED PWM}

The LED is connected to PB3, which corresponds to the OC2A output of
Timer2.

Timer2 is configured in Fast PWM mode with a prescaler of 64.

\begin{lstlisting}[language=C]
void LED_PWM_Init(void) {
    DDRB |= (1 << LED_PIN);

    TCCR2A = (1 << WGM20) |
             (1 << WGM21) |
             (1 << COM2A1);

    TCCR2B = (1 << CS22);
}
\end{lstlisting}

The LED brightness is controlled by changing the value of
\texttt{OCR2A}. The servo angle is mapped to the 8-bit PWM range from
0 to 255.

\begin{lstlisting}[language=C]
void LED_SetBrightnessByAngle(uint8_t angle) {
    if (angle > 180)
        angle = 180;

    uint8_t pwm_val =
        (uint16_t)angle * 255 / 180;

    OCR2A = pwm_val;
}
\end{lstlisting}

The relationship between the servo angle and the LED PWM value is:

\[
PWM_{\mathrm{LED}}
=
\frac{\theta}{180}\times255
\]

Thus:

\[
\theta = 0^\circ
\quad\Rightarrow\quad
PWM_{\mathrm{LED}} = 0
\]

and:

\[
\theta = 180^\circ
\quad\Rightarrow\quad
PWM_{\mathrm{LED}} = 255
\]

Therefore, the LED brightness increases when the servo angle increases.

\subsection{Servo and LED Operation}

The program changes the servo angle from 0 to 180 degrees in steps of
5 degrees. At the same time, the LED PWM value is changed according to
the servo angle.

After reaching 180 degrees, the servo angle is decreased back to
0 degrees.

\begin{lstlisting}[language=C]
for (uint8_t angle = 0; angle <= 180; angle += 5) {
    Servo_SetAngle(angle);
    LED_SetBrightnessByAngle(angle);

    _delay_ms(50);
}

for (int16_t angle = 180; angle >= 0; angle -= 5) {
    Servo_SetAngle((uint8_t)angle);
    LED_SetBrightnessByAngle((uint8_t)angle);

    _delay_ms(50);
}
\end{lstlisting}

The operating sequence is:

\[
0^\circ
\rightarrow
180^\circ
\rightarrow
0^\circ
\]

At the same time, the LED PWM changes from:

\[
0
\rightarrow
255
\rightarrow
0
\]

\subsection{Oscilloscope Measurement}

The PWM signal generated for the servo motor can be observed using an
oscilloscope.

The oscilloscope is used to observe:

\begin{itemize}
    \item PWM frequency
    \item PWM period
    \item Pulse width
    \item Change of pulse width when the servo angle changes
\end{itemize}

The expected PWM period is:

\[
T_{\mathrm{PWM}}
\approx
20\,\mathrm{ms}
\]

and the expected frequency is:

\[
f_{\mathrm{PWM}}
\approx
50\,\mathrm{Hz}
\]

The pulse width changes according to the servo angle and is expected
to be approximately within the range:

\[
0.5\,\mathrm{ms}
\leq
PW
\leq
2.5\,\mathrm{ms}
\]

When the servo angle changes, the oscilloscope should show that the
pulse width changes while the PWM period remains approximately
constant.

\subsection{Result}

Timer1 generates a PWM signal with a frequency of approximately 50 Hz
to control the servo motor. The servo position is controlled by
changing the PWM pulse width.

Timer2 generates PWM for the LED, and the LED brightness changes
according to the servo angle.

The oscilloscope can be used to verify the PWM period, frequency, and
pulse width of the servo control signal.


\section{Practice 5: Timer and Scheduler}

\subsection{Objective}

The objective of this practice is to use Timer0 to create a periodic
time base and implement a simple scheduler for tasks with different
execution periods.

The required tasks are:

\begin{itemize}
    \item Read ADC every 100 ms.
    \item Update LED PWM every 20 ms.
    \item Read sensor every 100 ms.
    \item Send data through UART every 1000 ms.
\end{itemize}

In this program, Timer0 is used to generate a 1 ms time base. Software
counters are then used to create the required task periods.

\subsection{Timer0 Initialization}

Timer0 is configured to generate a periodic interrupt. The timer uses
CTC mode with a prescaler of 64.

\begin{lstlisting}[language=C]
void Timer0_Init(void) {
    TCCR0A = (1 << WGM01);

    TCCR0B = (1 << CS01) |
             (1 << CS00);

    OCR0A = 249;

    TIMSK0 = (1 << OCIE0A);
}
\end{lstlisting}

With a CPU clock of 16 MHz and a prescaler of 64, the timer clock is:

\[
f_{\mathrm{Timer}}
=
\frac{16\,\mathrm{MHz}}{64}
=
250\,\mathrm{kHz}
\]

Therefore, each timer count has a duration of:

\[
T_{\mathrm{count}}
=
\frac{1}{250\,\mathrm{kHz}}
=
4\,\mu\mathrm{s}
\]

With:

\[
OCR0A=249
\]

the interrupt period is:

\[
T_{\mathrm{interrupt}}
=
(249+1)\times4\,\mu\mathrm{s}
=
1\,\mathrm{ms}
\]

Thus, Timer0 provides a 1 ms time base for the scheduler.

\subsection{Timer Interrupt}

The Timer0 Compare Match A interrupt sets a flag every 1 ms.

\begin{lstlisting}[language=C]
ISR(TIMER0_COMPA_vect) {
    timer_flag = 1;
}
\end{lstlisting}

The main loop checks this flag. When the flag is set, it is cleared
and the software task counters are increased.

\begin{lstlisting}[language=C]
if (timer_flag) {
    timer_flag = 0;

    t_10ms++;
    t_20ms++;
    t_100ms++;
    t_1000ms++;
}
\end{lstlisting}

The 1 ms timer therefore acts as the basic clock for all scheduled
tasks.

\subsection{10 ms Task}

The 10 ms counter is increased every 1 ms. When it reaches 10, the
counter is reset.

\begin{lstlisting}[language=C]
if (t_10ms >= 10) {
    t_10ms = 0;

    // 10 ms task
}
\end{lstlisting}

Therefore, the task is executed every:

\[
10\times1\,\mathrm{ms}
=
10\,\mathrm{ms}
\]

This provides a periodic task with a frequency of:

\[
f_{10ms}
=
\frac{1}{10\,\mathrm{ms}}
=
100\,\mathrm{Hz}
\]

\subsection{20 ms Task}

The 20 ms task is used to update the LED PWM value according to the
sensor value.

\begin{lstlisting}[language=C]
if (t_20ms >= 20) {
    t_20ms = 0;

    uint16_t temp_adc = adc_sensor_value;

    if (temp_adc < 300)
        temp_adc = 300;

    if (temp_adc > 555)
        temp_adc = 555;

    uint8_t pwm_out =
        (uint32_t)(temp_adc - 300) * 255 /
        (555 - 300);

    OCR0A = pwm_out;
}
\end{lstlisting}

The ADC value is limited to the range:

\[
300 \leq ADC \leq 555
\]

The value is then mapped to the 8-bit PWM range:

\[
0 \leq PWM \leq 255
\]

The mapping equation is:

\[
PWM
=
\frac{ADC-300}{555-300}
\times255
\]

The task is executed every:

\[
20\times1\,\mathrm{ms}
=
20\,\mathrm{ms}
\]

Therefore, the LED PWM update frequency is:

\[
f_{20ms}
=
\frac{1}{20\,\mathrm{ms}}
=
50\,\mathrm{Hz}
\]

\subsection{100 ms Task}

The sensor value is read from ADC channel 0 every 100 ms.

\begin{lstlisting}[language=C]
if (t_100ms >= 100) {
    t_100ms = 0;

    adc_sensor_value = ADC_Read(0);
}
\end{lstlisting}

The ADC reading is performed using the following function:

\begin{lstlisting}[language=C]
uint16_t ADC_Read(uint8_t channel) {
    ADMUX = (ADMUX & 0xF0) |
            (channel & 0x0F);

    ADCSRA |= (1 << ADSC);

    while (ADCSRA & (1 << ADSC));

    return ADC;
}
\end{lstlisting}

The task period is:

\[
100\times1\,\mathrm{ms}
=
100\,\mathrm{ms}
\]

Therefore, the sensor is read at a frequency of:

\[
f_{100ms}
=
\frac{1}{100\,\mathrm{ms}}
=
10\,\mathrm{Hz}
\]

\subsection{1000 ms Task}

The UART task is executed every 1000 ms. It sends the current ADC
value and PWM value to the computer.

\begin{lstlisting}[language=C]
if (t_1000ms >= 1000) {
    t_1000ms = 0;

    UART_SendString("Gia tri ADC Cam Bien: ");
    UART_SendNumber(adc_sensor_value);

    UART_SendString(" | Duty Cycle PWM: ");
    UART_SendNumber(OCR0A);

    UART_SendString("\r\n");
}
\end{lstlisting}

The task period is:

\[
1000\times1\,\mathrm{ms}
=
1000\,\mathrm{ms}
=
1\,\mathrm{s}
\]

Therefore, UART data is transmitted once per second:

\[
f_{1000ms}
=
\frac{1}{1\,\mathrm{s}}
=
1\,\mathrm{Hz}
\]

\subsection{Scheduler Structure}

The complete scheduler can be represented as a set of software
counters driven by the 1 ms timer interrupt:

\[
1\,\mathrm{ms}
\rightarrow
\begin{cases}
10\,\mathrm{ms}\rightarrow\text{Task 1}\\
20\,\mathrm{ms}\rightarrow\text{Task 2}\\
100\,\mathrm{ms}\rightarrow\text{Task 3}\\
1000\,\mathrm{ms}\rightarrow\text{Task 4}
\end{cases}
\]

The main loop checks each counter and executes the corresponding task
when its period has elapsed.

\subsection{Result}

Timer0 provides a 1 ms time base for the scheduler. Software counters
are used to generate different task periods without using a separate
timer for each task.

The resulting task periods are:

\begin{center}
\begin{tabular}{|c|c|}
\hline
\textbf{Task} & \textbf{Period} \\
\hline
Task 1 & 10 ms \\
\hline
Task 2 & 20 ms \\
\hline
Task 3 & 100 ms \\
\hline
Task 4 & 1000 ms \\
\hline
\end{tabular}
\end{center}

This scheduler allows tasks with different execution periods to run
from a common 1 ms time base.


\section{Practice 6: Control Step Motor}

\subsection{Objective}

The objective of this practice is to control a stepper motor using the
DRV8825 driver. The motor is controlled by generating pulse signals and
setting the rotation direction. The motor is commanded to move through
a sequence of predefined angles.

\subsection{DRV8825 Connections}

The control signals of the DRV8825 are connected to the ATmega328 as
follows:

\begin{itemize}
    \item PD2: EN (Enable)
    \item PD3: DIR (Direction)
    \item PD4: PUL (Pulse/Step)
\end{itemize}

The corresponding Arduino pins are defined as:

\begin{lstlisting}[language=C]
#define EN_PIN   2
#define DIR_PIN  3
#define PUL_PIN  4
\end{lstlisting}

The driver is enabled by setting the EN pin to LOW:

\begin{lstlisting}[language=C]
digitalWrite(EN_PIN, LOW);
\end{lstlisting}

\subsection{Step and Angle Configuration}

The motor is configured to use the following microstep setting:

\[
M0:M1:M2 = 1:0:0
\]

This configuration corresponds to half-step operation. In the tested
configuration, 400 pulses are required for one complete revolution of
360 degrees.

Therefore, the number of steps per degree is:

\[
StepsPerDegree
=
\frac{400}{360}
\]

The program defines this relationship as:

\begin{lstlisting}[language=C]
const float STEPS_PER_DEGREE = 400.0 / 360.0;
\end{lstlisting}

For a required angular displacement $\Delta\theta$, the number of
pulses is calculated by:

\[
N
=
\Delta\theta
\times
\frac{400}{360}
\]

where $N$ is the required number of pulses.

\subsection{Generating Step Pulses}

The \texttt{stepMotor()} function generates the required number of
pulses on the PUL pin.

\begin{lstlisting}[language=C]
void stepMotor(uint16_t steps, bool direction) {
    digitalWrite(DIR_PIN,
                 direction ? HIGH : LOW);

    for (uint16_t i = 0; i < steps; i++) {
        digitalWrite(PUL_PIN, HIGH);
        delayMicroseconds(1000);

        digitalWrite(PUL_PIN, LOW);
        delayMicroseconds(1000);
    }
}
\end{lstlisting}

The DIR pin determines the rotation direction. The PUL pin generates
the step pulses.

Each pulse consists of:

\[
1\,\mathrm{ms}\ \mathrm{HIGH}
+
1\,\mathrm{ms}\ \mathrm{LOW}
\]

Therefore, the pulse period is approximately:

\[
T_{\mathrm{pulse}}
=
2\,\mathrm{ms}
\]

and the corresponding pulse frequency is:

\[
f_{\mathrm{pulse}}
=
\frac{1}{2\,\mathrm{ms}}
=
500\,\mathrm{Hz}
\]

\subsection{Moving to a Target Angle}

The \texttt{moveToAngle()} function calculates the difference between
the target angle and the current angle.

\begin{lstlisting}[language=C]
void moveToAngle(uint16_t target_angle) {
    int16_t angle_diff =
        target_angle - current_angle;

    if (angle_diff == 0)
        return;

    bool direction =
        (angle_diff > 0);

    uint16_t abs_angle_diff =
        abs(angle_diff);

    uint16_t steps_to_move =
        round(abs_angle_diff *
              STEPS_PER_DEGREE);

    stepMotor(steps_to_move, direction);

    current_angle = target_angle;
}
\end{lstlisting}

The direction is determined from the sign of the angle difference:

\[
\Delta\theta
=
\theta_{\mathrm{target}}
-
\theta_{\mathrm{current}}
\]

If:

\[
\Delta\theta > 0
\]

the motor rotates in one direction.

If:

\[
\Delta\theta < 0
\]

the motor rotates in the opposite direction.

The absolute angular difference is then converted into the required
number of pulses.

\subsection{Target Angle Sequence}

The program defines the following target angles:

\begin{lstlisting}[language=C]
const uint16_t target_angles[] =
{
    0, 45, 90, 135, 180,
    225, 270, 315, 360
};
\end{lstlisting}

The motor therefore moves through the following sequence:

\[
0^\circ
\rightarrow
45^\circ
\rightarrow
90^\circ
\rightarrow
135^\circ
\rightarrow
180^\circ
\]

\[
\rightarrow
225^\circ
\rightarrow
270^\circ
\rightarrow
315^\circ
\rightarrow
360^\circ
\]

At each target angle, the program waits for one second:

\begin{lstlisting}[language=C]
delay(1000);
\end{lstlisting}

The motor therefore remains at each target position before moving to
the next position.

\subsection{Returning to the Initial Position}

After reaching the final target angle of 360 degrees, the program
resets the software position to 0 degrees:

\begin{lstlisting}[language=C]
current_angle = 0;
\end{lstlisting}

The loop then starts again from the first target angle.

Therefore, the complete operating sequence is:

\[
0^\circ
\rightarrow
45^\circ
\rightarrow
90^\circ
\rightarrow
135^\circ
\rightarrow
180^\circ
\]

\[
\rightarrow
225^\circ
\rightarrow
270^\circ
\rightarrow
315^\circ
\rightarrow
360^\circ
\rightarrow
0^\circ
\]

\subsection{Result}

The ATmega328 controls the stepper motor through the DRV8825 driver
using three control signals: Enable, Direction, and Pulse.

The motor position is controlled by the number of pulses generated on
the PUL pin. The DIR signal determines the rotation direction, while
the target angle is converted into the corresponding number of motor
steps.

The tested program successfully moves the motor through the predefined
angles from 0 to 360 degrees and repeats the sequence continuously.


\section{Practice 7: Input Capture and Timer1}

\subsection{Objective}

The objective of this practice is to use the Input Capture Unit of
Timer1 to measure the pulse width and period of an input signal.

The input signal is connected to the ICP1 input pin, which corresponds
to PB0 (D8) on the ATmega328.

\subsection{Input Capture Pin}

The ICP1 input is connected to PB0. The pin is configured as an input:

\begin{lstlisting}[language=C]
DDRB &= ~(1 << DDB0);
\end{lstlisting}

The Input Capture Unit uses the Timer1 counter to record the timer
value when an input edge is detected.

\subsection{Timer1 Configuration}

Timer1 is configured for Input Capture operation. The timer uses a
prescaler of 8.

\begin{lstlisting}[language=C]
void Timer1_ICP1_Init() {
    TCCR1A = 0;
    TCCR1B = 0;
    TCNT1 = 0;

    TCCR1B = (1 << ICNC1) |
             (1 << ICES1) |
             (1 << CS11);

    TIFR1 |= (1 << ICF1);
    TIMSK1 |= (1 << ICIE1);
}
\end{lstlisting}

With a CPU clock of 16 MHz and a prescaler of 8, the Timer1 clock is:

\[
f_{\mathrm{Timer1}}
=
\frac{16\,\mathrm{MHz}}{8}
=
2\,\mathrm{MHz}
\]

Therefore, one timer tick corresponds to:

\[
T_{\mathrm{tick}}
=
\frac{1}{2\,\mathrm{MHz}}
=
0.5\,\mu\mathrm{s}
\]

Thus, the Timer1 counter can be used to measure the time difference
between two captured edges.

\subsection{Input Capture Operation}

The Input Capture Unit generates an interrupt whenever the configured
input edge is detected.

The program initially detects the rising edge:

\begin{lstlisting}[language=C]
ICES1 = 1
\end{lstlisting}

When a rising edge is detected, the current Timer1 value is stored.
The program then changes the capture configuration to detect the
falling edge.

When the falling edge is detected, the Timer1 value is stored again.
The difference between the two values represents the duration of the
HIGH level of the input signal.

The capture process is:

\[
\text{Rising Edge}
\rightarrow
\text{Falling Edge}
\rightarrow
\text{Rising Edge}
\rightarrow
\cdots
\]

\subsection{Measuring Pulse Width}

The Input Capture interrupt handles the rising and falling edges.

\begin{lstlisting}[language=C]
ISR(TIMER1_CAPT_vect) {
    uint16_t current_icr = ICR1;

    if (edge_state == 0) {
        t_rising = current_icr;

        if (t_last_rising != 0) {
            period_ticks =
                t_rising - t_last_rising;
        }

        t_last_rising = t_rising;

        TCCR1B &= ~(1 << ICES1);
        edge_state = 1;
    }
    else {
        t_falling = current_icr;

        pulse_width_ticks =
            t_falling - t_rising;

        TCCR1B |= (1 << ICES1);
        edge_state = 0;

        capture_complete = true;
    }
}
\end{lstlisting}

For a rising edge followed by a falling edge:

\[
N_{\mathrm{width}}
=
T_{\mathrm{falling}}
-
T_{\mathrm{rising}}
\]

where $N_{\mathrm{width}}$ is the measured number of Timer1 ticks.

Since each tick corresponds to $0.5\,\mu\mathrm{s}$:

\[
T_{\mathrm{HIGH}}
=
N_{\mathrm{width}}
\times
0.5\,\mu\mathrm{s}
\]

\subsection{Measuring Signal Period}

The period is measured using two consecutive rising edges.

The program stores the previous rising-edge timestamp in
\texttt{t\_last\_rising}.

The period in Timer1 ticks is calculated as:

\[
N_{\mathrm{period}}
=
T_{\mathrm{rising,current}}
-
T_{\mathrm{rising,previous}}
\]

The corresponding period in microseconds is:

\[
T_{\mathrm{period}}
=
N_{\mathrm{period}}
\times
0.5\,\mu\mathrm{s}
\]

The program performs this calculation with:

\begin{lstlisting}[language=C]
period_ticks = t_rising - t_last_rising;
\end{lstlisting}

\subsection{Converting Timer Ticks to Time}

After a complete measurement is captured, the measured values are
converted from Timer1 ticks to microseconds.

\begin{lstlisting}[language=C]
uint32_t high_time_us = width_ticks / 2;
uint32_t period_us = period_t / 2;
\end{lstlisting}

Because one Timer1 tick corresponds to $0.5\,\mu\mathrm{s}$:

\[
T_{\mathrm{HIGH}}
=
\frac{N_{\mathrm{width}}}{2}
\;\mu\mathrm{s}
\]

and:

\[
T_{\mathrm{period}}
=
\frac{N_{\mathrm{period}}}{2}
\;\mu\mathrm{s}
\]

\subsection{UART Output}

The measured pulse width and period are transmitted through USART at
9600 baud.

\begin{lstlisting}[language=C]
USART_SendString("Do rong xung (HIGH): ");
USART_SendNum(high_time_us);

USART_SendString(" us | Chu ky: ");
USART_SendNum(period_us);

USART_SendString(" us\r\n");
\end{lstlisting}

The output therefore displays the measured HIGH pulse width and the
signal period in microseconds.

\subsection{Measurement Process}

The complete measurement process can be summarized as:

\[
\text{Input Signal}
\rightarrow
\text{ICP1}
\rightarrow
\text{Timer1 Capture}
\rightarrow
\text{Calculate Time}
\rightarrow
\text{UART}
\]

The rising edge is used to determine the beginning of the pulse and
to calculate the period. The falling edge is used to determine the
end of the HIGH pulse.

\subsection{Result}

Timer1 Input Capture is used to measure the timing characteristics of
the input signal on ICP1.

With a prescaler of 8 and a 16 MHz CPU clock, each Timer1 tick
corresponds to $0.5\,\mu\mathrm{s}$. The program uses this time
resolution to calculate the HIGH pulse width and signal period.

The measured values are displayed through the USART interface at
9600 baud.


\section{Practice 8: I2C Communication}

\subsection{Objective}

The objective of this practice is to understand the basic operation of
the I2C communication protocol and implement the main I2C control
functions using the hardware I2C module of the ATmega328.

The main I2C operations implemented in this practice are:

\begin{itemize}
    \item I2C initialization
    \item START condition
    \item Data transmission
    \item Data reception with ACK
    \item Data reception with NACK
    \item STOP condition
\end{itemize}

There is no external I2C device connected in this practice. Therefore,
the transmitted data is also displayed on LEDs for observation.

\subsection{I2C Initialization}

The ATmega328 hardware I2C peripheral is initialized using the
following function:

\begin{lstlisting}[language=C]
void i2c_init(void) {
    TWSR = 0x00;
    TWBR = 72;
    TWCR = (1 << TWEN);
}
\end{lstlisting}

The \texttt{TWEN} bit enables the I2C hardware module.

The value of \texttt{TWBR} is used to configure the I2C clock
frequency.

\subsection{I2C START Condition}

The START condition is generated using the following function:

\begin{lstlisting}[language=C]
void i2c_start(void) {
    TWCR = (1 << TWINT) |
           (1 << TWSTA) |
           (1 << TWEN);
}
\end{lstlisting}

The \texttt{TWSTA} bit requests the hardware to generate an I2C START
condition.

The normal I2C communication sequence begins with a START condition:

\[
\text{START}
\rightarrow
\text{Address}
\rightarrow
\text{Data}
\rightarrow
\text{STOP}
\]

\subsection{I2C Data Transmission}

The \texttt{i2c\_write()} function sends one byte through the I2C
interface.

\begin{lstlisting}[language=C]
void i2c_write(uint8_t data) {
    PORTB = (PORTB & 0xFC) |
            (data & 0x03);

    PORTD = (PORTD & 0x03) |
            (data & 0xFC);

    TWDR = data;

    TWCR = (1 << TWINT) |
           (1 << TWEN);
}
\end{lstlisting}

Before transmitting the byte through the I2C peripheral, the data is
also mapped to the LED outputs.

The two least significant bits are displayed on PB0 and PB1:

\[
D_1D_0
\rightarrow
PB1,PB0
\]

The remaining six bits are displayed on PD2 through PD7:

\[
D_7D_6D_5D_4D_3D_2
\rightarrow
PD7,PD6,PD5,PD4,PD3,PD2
\]

Therefore, the complete 8-bit transmitted value can be observed using
the LEDs.

\subsection{I2C Address Transmission}

The sensor address used in the test sequence is:

\[
0x68
\]

For an I2C write operation, the 7-bit address is shifted left by one
bit and the write bit is added:

\[
Address
=
(0x68 << 1) | 0
\]

This value is transmitted using:

\begin{lstlisting}[language=C]
i2c_write((0x68 << 1) | 0);
\end{lstlisting}

The resulting I2C address byte contains both the sensor address and the
write indication.

\subsection{I2C Data Reception}

Two functions are provided for receiving data.

The first function requests an ACK after receiving a byte:

\begin{lstlisting}[language=C]
uint8_t i2c_read_ack(void) {
    TWCR = (1 << TWINT) |
           (1 << TWEN) |
           (1 << TWEA);

    return TWDR;
}
\end{lstlisting}

The second function receives a byte without sending an ACK:

\begin{lstlisting}[language=C]
uint8_t i2c_read_nack(void) {
    TWCR = (1 << TWINT) |
           (1 << TWEN);

    return TWDR;
}
\end{lstlisting}

The ACK and NACK functions are included to implement the two basic
types of I2C read operation.

In this practice, there is no external I2C device providing data, so
these functions are not used in the main test sequence.

\subsection{I2C STOP Condition}

The communication is terminated using the STOP condition:

\begin{lstlisting}[language=C]
void i2c_stop(void) {
    TWCR = (1 << TWINT) |
           (1 << TWSTO) |
           (1 << TWEN);
}
\end{lstlisting}

The \texttt{TWSTO} bit requests the hardware to generate an I2C STOP
condition.

Therefore, a complete transmission has the following structure:

\[
\text{START}
\rightarrow
\text{I2C Address}
\rightarrow
\text{Data}
\rightarrow
\text{STOP}
\]

\subsection{Test Data Sequence}

The program uses a sequence of bytes representing typical register
addresses and configuration data of an I2C sensor:

\begin{lstlisting}[language=C]
uint8_t sample_data[] = {
    0x75,
    0x6B,
    0x00,
    0x3B,
    0x1C,
    0x08,
    0x43,
    0x1A
};
\end{lstlisting}

These values represent example commands such as:

\begin{itemize}
    \item \texttt{0x75}: WHO\_AM\_I register
    \item \texttt{0x6B}: PWR\_MGMT\_1 register
    \item \texttt{0x00}: Wake-up configuration data
    \item \texttt{0x3B}: Accelerometer data register
    \item \texttt{0x1C}: Accelerometer configuration register
    \item \texttt{0x08}: Accelerometer range configuration
    \item \texttt{0x43}: Gyroscope data register
    \item \texttt{0x1A}: Configuration register
\end{itemize}

The program sends one byte from this sequence during each iteration.

\subsection{Communication Sequence}

The main loop performs the following operations:

\begin{lstlisting}[language=C]
while (1) {
    i2c_start();

    i2c_write((0x68 << 1) | 0);

    i2c_write(sample_data[current_index]);

    i2c_stop();

    current_index =
        (current_index + 1) % total_patterns;

    _delay_ms(500);
}
\end{lstlisting}

The resulting communication sequence is:

\[
\text{START}
\rightarrow
\text{Address + Write}
\rightarrow
\text{Register/Data}
\rightarrow
\text{STOP}
\]

A new byte from the test sequence is transmitted every 500 ms.

\subsection{LED Data Display}

The transmitted byte is simultaneously displayed on the LEDs connected
to PB0, PB1, and PD2 through PD7.

Therefore, the 8-bit data can be visually observed as:

\[
D_7D_6D_5D_4D_3D_2D_1D_0
\]

with the corresponding LED mapping:

\[
\begin{array}{c|c}
\text{Data bit} & \text{LED output} \\
\hline
D_0 & PB0 \\
D_1 & PB1 \\
D_2 & PD2 \\
D_3 & PD3 \\
D_4 & PD4 \\
D_5 & PD5 \\
D_6 & PD6 \\
D_7 & PD7
\end{array}
\]

This allows the transmitted byte to be observed without requiring an
external I2C device.

\subsection{Result}

The practice implements the basic I2C operations using the ATmega328
hardware I2C peripheral.

The program generates START and STOP conditions, sends the I2C address
and data bytes, and provides functions for receiving data with ACK or
NACK.

Because no external I2C device is connected, the transmitted bytes are
displayed on LEDs for observation. The test sequence sends one byte
every 500 ms.


\section{Practice 9: Serial Control of a Servo Motor Using ATmega328}

\subsection{Objective}

The objective of this practice is to control a servo motor through
serial communication with a computer.

The program receives an angle command through the serial interface,
checks whether the command is valid, and moves the servo to the
requested angle if the value is within the permitted range.

The valid servo angle range is:

\[
0^\circ \leq \theta \leq 180^\circ
\]

\subsection{Servo Configuration}

The servo motor is connected to digital pin 9.

\begin{lstlisting}[language=C]
const int SERVO_PIN = 9;
const int MIN_ANGLE = 0;
const int MAX_ANGLE = 180;
\end{lstlisting}

The servo is initialized in the \texttt{setup()} function:

\begin{lstlisting}[language=C]
Serial.begin(9600);

servo.attach(SERVO_PIN);
servo.write(90);
\end{lstlisting}

The serial communication is configured at:

\[
Baudrate = 9600\ \mathrm{baud}
\]

The servo is initially positioned at:

\[
\theta = 90^\circ
\]

\subsection{Receiving Serial Commands}

The program continuously checks whether data is available through the
serial interface.

\begin{lstlisting}[language=C]
while (Serial.available() > 0)
{
    char receivedChar = Serial.read();

    if (receivedChar == '\n' ||
        receivedChar == '\r')
    {
        if (inputBuffer.length() > 0)
        {
            processCommand(inputBuffer);
            inputBuffer = "";
        }
    }
    else
    {
        inputBuffer += receivedChar;
    }
}
\end{lstlisting}

The received characters are stored in \texttt{inputBuffer}. When a
newline or carriage return is received, the accumulated string is
processed as one command.

\subsection{Command Validation}

Before converting the command into an angle, the program checks
whether the input contains a valid numerical value.

\begin{lstlisting}[language=C]
bool validNumber = true;

for (unsigned int i = 0;
     i < command.length(); i++)
{
    if (!isDigit(command[i]) &&
        !(i == 0 && command[i] == '-'))
    {
        validNumber = false;
        break;
    }
}
\end{lstlisting}

If the command contains an invalid character, the program rejects the
command and sends an error message:

\begin{lstlisting}[language=C]
if (!validNumber)
{
    Serial.print("\nERROR: Invalid command.");
    return;
}
\end{lstlisting}

\subsection{Checking the Servo Angle}

After the command has been verified as a number, it is converted into
an integer angle:

\begin{lstlisting}[language=C]
int angle = command.toInt();
\end{lstlisting}

The program then checks whether the angle is within the permitted
range:

\[
0^\circ \leq \theta \leq 180^\circ
\]

\begin{lstlisting}[language=C]
if (angle < MIN_ANGLE ||
    angle > MAX_ANGLE)
{
    Serial.print(
        "\nERROR: Angle must be between "
        "0 and 180 degrees.\n"
    );
    return;
}
\end{lstlisting}

If the angle is outside this range, the servo is not moved and an error
message is sent to the computer.

\subsection{Controlling the Servo}

If the command is valid and the angle is within the permitted range,
the servo is moved to the requested position:

\begin{lstlisting}[language=C]
servo.write(angle);
\end{lstlisting}

The program then sends a confirmation message:

\begin{lstlisting}[language=C]
Serial.print("\nOK: Servo moved to ");
Serial.print(angle);
Serial.println(" degrees.");
\end{lstlisting}

Therefore, a valid command follows the sequence:

\[
\text{Serial Command}
\rightarrow
\text{Validate}
\rightarrow
\text{Check Range}
\rightarrow
\text{Move Servo}
\rightarrow
\text{Confirmation}
\]

An invalid command follows:

\[
\text{Serial Command}
\rightarrow
\text{Validate}
\rightarrow
\text{Reject}
\rightarrow
\text{Error Message}
\]

\subsection{Example of Operation}

For a valid command such as:

\begin{lstlisting}
90
\end{lstlisting}

the program moves the servo to 90 degrees and sends a confirmation
message.

For a command outside the permitted range, such as:

\begin{lstlisting}
200
\end{lstlisting}

the program rejects the command and sends an error message indicating
that the angle must be between 0 and 180 degrees.

For a command containing invalid characters, the program reports:

\begin{lstlisting}
ERROR: Invalid command.
\end{lstlisting}

\subsection{Result}

The ATmega328 receives servo angle commands from a computer through
the serial interface at 9600 baud.

The program validates the received command before controlling the
servo. Valid angles from 0 to 180 degrees are accepted and used to
control the servo, while invalid commands or out-of-range angles are
rejected with an error message.

After each valid command, the program sends a confirmation containing
the commanded servo angle.


\section{Practice 10: Motion Detection Using PIR Sensor}

\subsection{Objective}

The objective of this practice is to detect motion using a PIR sensor
connected to the ATmega328.

When the PIR sensor detects motion, the LED is turned on and a status
message is sent through UART. When motion is no longer detected, the
LED is turned off and the corresponding status message is sent.

The program also prevents repeated UART messages when the PIR state
does not change.

\subsection{Hardware Configuration}

The PIR sensor is connected to digital pin 2 and the LED is connected
to digital pin 13.

\begin{lstlisting}[language=C]
const int PIR_PIN = 2;
const int LED_PIN = 13;
\end{lstlisting}

The PIR pin is configured as an input and the LED pin is configured as
an output:

\begin{lstlisting}[language=C]
pinMode(PIR_PIN, INPUT);
pinMode(LED_PIN, OUTPUT);
\end{lstlisting}

The serial interface is configured at:

\[
Baudrate = 9600\ \mathrm{baud}
\]

\subsection{Reading the PIR Sensor}

The PIR sensor state is read using \texttt{digitalRead()}:

\begin{lstlisting}[language=C]
int currentState = digitalRead(PIR_PIN);
\end{lstlisting}

The sensor is configured as active-low. Therefore, a LOW signal
indicates motion detection in this implementation.

The states are interpreted as:

\[
PIR =
\begin{cases}
LOW & \text{Motion detected}\\
HIGH & \text{No motion}
\end{cases}
\]

\subsection{Detecting a State Change}

The program stores the previous PIR state in
\texttt{lastState}:

\begin{lstlisting}[language=C]
int lastState = LOW;
\end{lstlisting}

The current state is compared with the previous state:

\begin{lstlisting}[language=C]
if (currentState != lastState)
{
    ...
}
\end{lstlisting}

Therefore, the motion status is processed only when the PIR state
changes.

This prevents the program from repeatedly sending the same UART
message while the sensor remains in the same state.

\subsection{Motion Detected}

When the PIR output changes to LOW, the program determines that motion
has been detected.

\begin{lstlisting}[language=C]
if (currentState == LOW)
{
    digitalWrite(LED_PIN, HIGH);

    Serial.println("Motion detected.");
}
\end{lstlisting}

The LED is turned on and the following message is sent through UART:

\begin{lstlisting}
Motion detected.
\end{lstlisting}

\subsection{No Motion}

When the PIR output changes to HIGH, the program determines that
motion is no longer detected.

\begin{lstlisting}[language=C]
else
{
    digitalWrite(LED_PIN, LOW);

    Serial.println("No motion.");
}
\end{lstlisting}

The LED is turned off and the following message is sent through UART:

\begin{lstlisting}
No motion.
\end{lstlisting}

\subsection{Updating the Previous State}

After processing the state change, the current PIR state is stored as
the previous state:

\begin{lstlisting}[language=C]
lastState = currentState;
\end{lstlisting}

The next loop iteration then compares the new sensor state with this
stored value.

The complete state-change process is:

\[
\text{Read PIR}
\rightarrow
\text{Compare with Previous State}
\rightarrow
\text{State Changed?}
\]

If the state has changed:

\[
\text{Process State}
\rightarrow
\text{Update LED}
\rightarrow
\text{Send UART Message}
\]

If the state has not changed, no new message is transmitted.

\subsection{Result}

The ATmega328 successfully reads the digital output of the PIR sensor
and determines the motion state.

When motion is detected, the LED is turned on and the message
``Motion detected.'' is sent through UART. When there is no motion, the
LED is turned off and the message ``No motion.'' is sent.

The program only sends a new status message when the PIR state changes,
preventing repeated transmission of the same status.


\section{Practice 11: LDR and Light Detection Using ATmega328}

\subsection{Objective}

The objective of this practice is to measure light intensity using an
LDR connected to the ADC of the ATmega328.

The ADC value is used for three purposes:

\begin{itemize}
    \item Determine whether the environment is dark or bright.
    \item Control an LED according to the light condition.
    \item Display the ADC value on a 4-digit 7-segment display.
\end{itemize}

The ADC value and light condition are also transmitted through UART.

\subsection{Hardware Configuration}

The LDR is connected to analog input A0:

\begin{lstlisting}[language=C]
const int LIGHT_PIN = A0;
\end{lstlisting}

The LED is connected to digital pin 13:

\begin{lstlisting}[language=C]
const int LED_PIN = 13;
\end{lstlisting}

The 4-digit 7-segment display uses the following segment pins:

\begin{lstlisting}[language=C]
const int SEGMENT_PINS[8] = {
    3, 4, 5, 6, 7, 8, 9, 10
};
\end{lstlisting}

The four digit-selection pins are:

\begin{lstlisting}[language=C]
const int DIGIT_PINS[4] = {
    11, 12, A1, A2
};
\end{lstlisting}

The 7-segment display uses active-low logic. Therefore:

\[
LOW = ON
\]

and:

\[
HIGH = OFF
\]

\subsection{Reading the LDR}

The analog value from the LDR is continuously read using the
ATmega328 ADC:

\begin{lstlisting}[language=C]
int adcValue = analogRead(LIGHT_PIN);
\end{lstlisting}

The ADC value is represented by:

\[
0 \leq ADC \leq 1023
\]

The current ADC value is used directly for light-condition detection,
UART transmission, and display.

\subsection{Dark and Bright Detection}

A threshold value is defined to distinguish between dark and bright
conditions:

\begin{lstlisting}[language=C]
const int DARK_THRESHOLD = 400;
\end{lstlisting}

The light condition is determined using:

\[
Condition =
\begin{cases}
\text{Dark}, & ADC < 400\\
\text{Bright}, & ADC \geq 400
\end{cases}
\]

The program controls the LED according to this condition:

\begin{lstlisting}[language=C]
if (adcValue < DARK_THRESHOLD)
{
    digitalWrite(LED_PIN, HIGH);
}
else
{
    digitalWrite(LED_PIN, LOW);
}
\end{lstlisting}

Therefore, the LED is turned on when the measured ADC value is below
the dark threshold and turned off when the value is equal to or above
the threshold.

\subsection{Displaying the ADC Value}

The ADC value is displayed using a 4-digit 7-segment display.

The program separates the ADC value into four decimal digits:

\begin{lstlisting}[language=C]
digits[0] = (number / 1000) % 10;
digits[1] = (number / 100) % 10;
digits[2] = (number / 10) % 10;
digits[3] = number % 10;
\end{lstlisting}

For example, if:

\[
ADC = 523
\]

the four digits are:

\[
0,\ 5,\ 2,\ 3
\]

The display therefore represents:

\[
0523
\]

\subsection{7-Segment Digit Patterns}

The program stores the segment patterns for the digits from 0 to 9:

\begin{lstlisting}[language=C]
const byte DIGIT_PATTERN[10] = {
    0b00111111,
    0b00000110,
    0b01011011,
    0b01001111,
    0b01100110,
    0b01101101,
    0b01111101,
    0b00000111,
    0b01111111,
    0b01101111
};
\end{lstlisting}

The \texttt{setSegments()} function converts a digit pattern into the
corresponding segment outputs.

\begin{lstlisting}[language=C]
void setSegments(byte pattern)
{
    for (int i = 0; i < 8; i++)
    {
        if (bitRead(pattern, i))
        {
            digitalWrite(SEGMENT_PINS[i], LOW);
        }
        else
        {
            digitalWrite(SEGMENT_PINS[i], HIGH);
        }
    }
}
\end{lstlisting}

Because the display is active-low, a LOW output turns a segment on.

\subsection{4-Digit Multiplexing}

The four digits are controlled using multiplexing. Only one digit is
activated at a time, while the program rapidly switches between the
four digits.

The \texttt{displayNumber()} function performs this operation:

\begin{lstlisting}[language=C]
for (int i = 0; i < 4; i++)
{
    turnOffDigits();

    setSegments(DIGIT_PATTERN[digits[i]]);

    digitalWrite(DIGIT_PINS[i], LOW);

    delay(2);

    digitalWrite(DIGIT_PINS[i], HIGH);
}
\end{lstlisting}

The sequence is:

\[
\text{Turn OFF all digits}
\rightarrow
\text{Set segments}
\rightarrow
\text{Turn ON one digit}
\rightarrow
\text{Short delay}
\rightarrow
\text{Turn OFF digit}
\]

This process is repeated continuously to maintain the appearance of a
four-digit display.

\subsection{Updating the Display Value}

The ADC is read continuously, but the value used for the 7-segment
display is updated every 100 ms.

\begin{lstlisting}[language=C]
static int displayValue = 0;
static unsigned long lastDisplayUpdate = 0;

if (millis() - lastDisplayUpdate >= 100)
{
    lastDisplayUpdate = millis();

    displayValue = adcValue;
}
\end{lstlisting}

Therefore, the ADC measurement and the display update operate at
different rates:

\[
\text{ADC}
\rightarrow
\text{Continuous Reading}
\]

while:

\[
\text{7-Segment Value}
\rightarrow
\text{Update Every 100 ms}
\]

\subsection{UART Output}

The program sends the ADC value and the corresponding light condition
through UART.

The UART interface is configured at:

\[
Baudrate = 9600\ \mathrm{baud}
\]

The output is updated every 100 ms:

\begin{lstlisting}[language=C]
if (millis() - lastSerialUpdate >= 100)
{
    lastSerialUpdate = millis();

    Serial.print("ADC: ");
    Serial.print(adcValue);

    if (adcValue < DARK_THRESHOLD)
    {
        Serial.println(" - Dark");
    }
    else
    {
        Serial.println(" - Bright");
    }
}
\end{lstlisting}

The transmitted information has the following form:

\begin{lstlisting}
ADC: 350 - Dark
\end{lstlisting}

or:

\begin{lstlisting}
ADC: 700 - Bright
\end{lstlisting}

depending on the measured light level.

\subsection{Overall Operation}

The overall operation of the program can be summarized as:

\[
\text{LDR}
\rightarrow
\text{ADC}
\rightarrow
\text{ADC Value}
\]

The ADC value is then used by three independent functions:

\[
\text{ADC Value}
\rightarrow
\begin{cases}
\text{Dark/Bright Detection}\\
\text{LED Control}\\
\text{7-Segment Display}\\
\text{UART Output}
\end{cases}
\]

The 7-segment display shows the latest ADC value, while the LED
indicates the detected light condition.

\subsection{Result}

The ATmega328 continuously reads the analog signal from the LDR using
its ADC.

When the ADC value is below 400, the environment is classified as
dark and the LED is turned on. When the ADC value is 400 or higher,
the environment is classified as bright and the LED is turned off.

The ADC value is displayed on a 4-digit active-low 7-segment display
using multiplexing. The displayed value is updated every 100 ms.

The ADC value and the corresponding light condition are also
transmitted through UART at 9600 baud.

\end{document}